\chapter{Things Every User Needs to Know}

This chapter explains the handful of ideas that appear on nearly every screen.
Read it once and the rest of the manual will make sense.

\section{The scan box}

Wherever the system needs to know which garments you are holding, you will see a
large box like the one in Figure~\ref{fig:scanbox}. This is the \textbf{scan box}.

\shotsmall{71-transfers-receive}{The scan box. The big number on the blue strip
counts the tags you have read so far.}{fig:scanbox}

There are three ways tags get into that box.

\begin{steps}
  \item \textbf{A handheld reader.} Most hand scanners are set up to behave like
        a keyboard. Click inside the box once, then pull the trigger. The tag
        numbers type themselves in, one per line.
  \item \textbf{A fixed reader.} At a portal or tunnel the reader sends the tags
        to the system by itself. You do not have to do anything.
  \item \textbf{The \btn{Simulate reader} button.} This is for training and
        testing. It fills the box with the tags that really are in that location,
        so you can practise without hardware.
\end{steps}

\begin{tipbox}
Reading the same tag twice does no harm. The system counts each garment once,
however many times its tag is seen. If in doubt, scan again.
\end{tipbox}

The blue strip shows how many tags are ready. \btn{Clear} empties the box and
starts again.

\section{Coloured labels}

Small coloured labels appear all over the system. The colour tells you the
meaning at a glance.

\begin{table}[H]
\centering
\small
\begin{tabular}{@{}p{26mm}p{26mm}p{86mm}@{}}
\toprule
\textbf{Colour} & \textbf{Examples} & \textbf{Meaning} \\
\midrule
Green   & \texttt{PASS}, \texttt{RECEIVED}, \texttt{READY}
        & Everything is fine. Nothing to do. \\
Blue    & \texttt{IN TRANSIT}, \texttt{DISPATCHED}
        & Normal, in progress. Somebody else has the next step. \\
Amber   & \texttt{VARIANCE}, \texttt{REWORK}, \texttt{HOLD}
        & Needs a person to look at it. \\
Red     & \texttt{FAIL}, \texttt{SCRAP}, \texttt{MISSING}
        & Something is wrong, or the garment has been rejected. \\
Grey    & \texttt{PENDING}, \texttt{CUT}
        & Waiting its turn. Normal. \\
\bottomrule
\end{tabular}
\caption{What the colours mean.}
\end{table}

\section{Words you will see}

\begin{table}[H]
\centering
\small
\begin{tabular}{@{}p{34mm}p{104mm}@{}}
\toprule
\textbf{Word} & \textbf{What it means in plain language} \\
\midrule
Tag / EPC &
  The electronic label on the garment, and the unique number stored inside it.
  \texttt{EPC} is just the technical name for that number. \\
Article &
  One single garment. Every article has its own tag and its own serial number. \\
Serial number &
  A readable name for a garment, like \typed{ART-260816-000123}. Easier to say
  aloud than the tag number. \\
Bundle &
  A tied stack of cut pieces --- for example fifty size 32 fronts and backs ---
  handed from cutting to the sewing line. \\
Transfer note &
  The document created when one department sends goods to another. It lists
  exactly what was sent. \\
Batch &
  A group of garments moved together, usually all the same design and size. \\
Variance &
  The count on arrival did not match the count that was sent. \\
WIP &
  ``Work in progress'': everything currently on the factory floor, not yet shipped. \\
Commissioning &
  The moment a tag is attached to a garment and registered in the system.
  After this, the garment exists as far as the computer is concerned. \\
\bottomrule
\end{tabular}
\caption{The vocabulary of the system.}
\end{table}

\section{Tables: sorting, searching and opening}

Most screens show a table. Three things are always true:

\begin{bullets}
  \item \textbf{Rows can be clicked.} Clicking a row opens the full detail of that item.
  \item \textbf{Counts are on the right.} Numbers line up so you can compare them down a column.
  \item \textbf{\btn{Export CSV} saves the table} to a file you can open in Excel.
\end{bullets}

\section{If something goes wrong}

The system tells you when it cannot do what you asked, in a small panel at the
bottom right of the screen. It always explains why, in words --- for example
``Currently in Sorting Station, not Washing''.

\begin{warnbox}
The system refuses things on purpose to protect the count. If a garment is
rejected, it is nearly always because it is genuinely somewhere else, or its tag
belongs to another garment. Put that piece aside and tell your supervisor rather
than trying again repeatedly.
\end{warnbox}

Chapter~\ref{ch:trouble} lists the messages you are most likely to meet and what
to do about each one.

\section{Moving goods between departments}

Every station does these two things, whatever else its job is. Read this once;
the rest of the manual refers back to it.

\section{Receiving goods into your section}
\label{sec:receive}

Open \menu{Transfers} from the left menu. The top table, \emph{Waiting to be
received}, lists everything other departments have sent you.

\shot{70-transfers-inbox}{The Transfers screen. Everything sent to your section
is at the top, with how long it has been waiting.}{fig:op-inbox}

The \field{Waiting} column turns amber after four hours and red after eight.
That is your cue to deal with it.

\begin{steps}
  \item Find the row for the trolley in front of you. The \field{Document}
        number is printed on the note that came with it.
  \item Press \btn{Receive} on that row.
  \item Bulk-read everything that arrived into the scan box.
  \item Press \btn{Receive and tally}.
\end{steps}

\subsection{Reading the result}

The system now compares what arrived against what was sent, and shows five boxes.

\shotsmall{72-transfers-tally}{The result of a count. Here two pieces were not
found, so the handover is marked as a variance and the missing garments are
listed by name.}{fig:op-tally}

\begin{table}[H]
\centering
\small
\begin{tabular}{@{}p{24mm}p{110mm}@{}}
\toprule
\textbf{Box} & \textbf{Meaning} \\
\midrule
Expected & How many the sending department said they sent. \\
Received & How many you have just found. \\
Missing  & On the note but not found. \\
Extra    & Found but not on the note. \\
Result   & \texttt{MATCHED} if everything agrees, \texttt{VARIANCE} if not. \\
\bottomrule
\end{tabular}
\caption{The five boxes shown after a count.}
\end{table}

\begin{tipbox}
\textbf{If it says VARIANCE, do not panic.} Look through the trolley again and
scan a second time. Anything found is added to the count and the variance clears
by itself. Only if the pieces genuinely cannot be found does a supervisor need to
close the variance --- your supervisor has to close it.
\end{tipbox}

\section{Sending goods to the next department}
\label{sec:dispatch}

\begin{steps}
  \item On the \menu{Transfers} screen press \btn{+ Dispatch a batch}.
  \item Choose the department in \field{Send to}.
  \item Type a \field{Batch reference} if your section uses one, for example
        \typed{WASH-LOT-07}. This is printed on the note so the receiving
        department can find it.
  \item Bulk-read everything that is leaving.
  \item Press \btn{Check what was read} if you want to see a breakdown first.
  \item Press \btn{Generate transfer note}.
\end{steps}

\shotsmall{73-transfers-dispatch}{Sending a batch onward. The scan box has been
filled by the reader; the note is created when you press the blue button.}{fig:op-dispatch}

The printable note opens in a new tab automatically. Print it and send it with
the trolley.

\begin{warnbox}
If some tags are refused you will see a list with a reason against each one, such
as ``Currently in Finishing, not Sorting Station''. Those garments are not on the
note. Put them aside --- they are in the wrong place and need a supervisor.
\end{warnbox}

\section{Looking back at what your section has moved}

Below the inbox on the \menu{Transfers} screen is \emph{Transfer history for this
section}: every handover your department has been part of, in both directions.

\shot{74-transfers-history}{Transfer history. \texttt{IN} rows are things you
received, \texttt{OUT} rows are things you sent. The \field{Sent} and
\field{Received} columns should match.}{fig:op-history}

\begin{table}[H]
\centering
\small
\begin{tabular}{@{}p{24mm}p{110mm}@{}}
\toprule
\textbf{Column} & \textbf{What it tells you} \\
\midrule
Document & The transfer note number, printed on the paper that travelled with the goods. \\
Route    & Which department sent it and which received it. \\
Batch    & Your own reference for the load, if one was typed in. \\
Sent \newline Received & The two counts. If they differ, the difference is in the next two columns. \\
Missing \newline Extra & How many were on the note but not found, and found but not on the note. \\
Status   & \texttt{RECEIVED} when it is settled, \texttt{VARIANCE} while a difference is still open. \\
\bottomrule
\end{tabular}
\caption{Reading a row of transfer history.}
\end{table}

\begin{tipbox}
Use the \field{Section} box at the top right to look at a different department.
This is the quickest way to answer ``did you ever send us that batch?'' --- the
answer is a document number, a time and a count.
\end{tipbox}
